\begin{justify}

\textbf{DEFINICIÓN DE FUNCIÓN CUADRÁTICA}

La función cuadrática es una función polinómica de segundo grado que tiene la forma:
\[
f(x) = ax^2 + bx + c
\]
donde $a \neq 0$.

\textbf{PASOS PARA ELABORAR EL GRÁFICO DE UNA FUNCIÓN CUADRÁTICA}

\begin{enumerate}

\item \textbf{Identificar coeficientes:}  
Determina los valores de $a$, $b$ y $c$ en la función.

\item \textbf{Dirección de la parábola:}  
Si $a > 0$, la parábola abre hacia arriba.  
Si $a < 0$, la parábola abre hacia abajo.

\item \textbf{Vértice:}  
Se calcula con la fórmula:
\[
x = \frac{-b}{2a}
\]

\item \textbf{Intersección con el eje Y:}  
Es el valor de $c$, es decir el punto $(0, c)$.

\item \textbf{Graficar puntos:}  
Sustituye valores de $x$ para obtener puntos y dibujar la parábola.

\end{enumerate}

\end{justify}
